% =====================================================================
% Requiere: pgfgantt, xcolor[table], tabularx, booktabs, float
% Los colores se declaran con \providecolor: si ya los tienes definidos
% en tu preámbulo se respetan los tuyos.
% =====================================================================
\providecolor{azulPastelApagado}{HTML}{5B7FA6}
\providecolor{azulclaro}{HTML}{DCE6F2}
\providecolor{textoOscuro}{HTML}{22303C}
\providecolor{grisPlan}{HTML}{D5D9DE}
\providecolor{verdeOk}{HTML}{5FA777}
\providecolor{ambarParcial}{HTML}{E5A93D}
\providecolor{rojoFallo}{HTML}{D1584F}
\providecolor{azulNuevo}{HTML}{4A7FBF}

% Estilo común de los dos diagramas
\ganttset{cmp/.style={
    hgrid, vgrid,
    x unit=0.6cm,
    y unit title=0.65cm,
    y unit chart=0.62cm,
    title height=1,
    title/.style={fill=azulPastelApagado, draw=azulPastelApagado!70!black},
    title label font=\bfseries\small\color{white},
    title label node/.append style={text=white},
    bar label font=\footnotesize\color{textoOscuro},
    group label font=\footnotesize\bfseries\color{textoOscuro},
    group/.append style={fill=textoOscuro, draw=textoOscuro},
    group right peak width=0.2,
    group left peak width=0.2,
    vrule/.append style={rojoFallo, line width=1.2pt},
    vrule label font=\scriptsize\bfseries\color{rojoFallo},
    milestone/.append style={fill=rojoFallo, draw=rojoFallo!70!black}
}}

% \planreal{etiqueta}{ini plan}{fin plan}{ini real}{fin real}{color}
% Barra gris ancha = planificado; barra fina de color = realidad
\newcommand{\planreal}[6]{%
  \ganttbar[bar height=0.85, bar/.append style={fill=grisPlan, draw=grisPlan!60!black}]{#1}{#2}{#3}%
  \ganttbar[bar height=0.35, bar/.append style={fill=#6, draw=#6!60!black}]{}{#4}{#5}%
  \ganttnewline}

% Leyenda reutilizable
\newcommand{\cajita}[1]{\fcolorbox{black!40}{#1}{\rule{0pt}{0.8ex}\rule{1.1em}{0pt}}}

\section{Análisis Visual de las Iteraciones}

Esta sección resume de forma visual qué se planificó, qué se consiguió realmente y cómo se traslada la deuda técnica a la siguiente iteración.

\subsection{Iteración 1: Planificado frente a Realidad}

En cada fila, la \textbf{barra gris ancha} muestra lo previsto en el plan inicial y la \textbf{barra fina de color} refleja lo ocurrido. La línea roja vertical marca la baja no prevista de un integrante del equipo de backend.

\begin{center}
\small
\cajita{grisPlan}~Planificado \quad
\cajita{verdeOk}~Completado \quad
\cajita{ambarParcial}~Parcial / con retraso \quad
\cajita{rojoFallo}~No completado (pasa a Iteración 2)
\end{center}

\begin{figure}[H]
\centering
\resizebox{\textwidth}{!}{%
\begin{ganttchart}[cmp]{1}{21}
  \gantttitle{Octubre 2026}{10}
  \gantttitle{Noviembre 2026}{11} \\
  \gantttitlelist{22,23,24,25,26,27,28,29,30,31,1,2,3,4,5,6,7,8,9,10,11}{1} \\

  \ganttgroup{Front-End y Accesibilidad}{1}{21} \\
  \planreal{Bocetos de todas las interfaces}{1}{8}{1}{8}{verdeOk}
  \planreal{Implementar inicios de sesión}{9}{11}{9}{16}{ambarParcial}
  \planreal{Implementar vistas (actividad, agenda, asignación)}{12}{18}{12}{21}{ambarParcial}
  \planreal{Integración Front--Back}{19}{20}{19}{21}{rojoFallo}
  \planreal{Revisión del contenido generado}{21}{21}{21}{21}{rojoFallo}

  \ganttgroup{Back-End e Infraestructura}{1}{21} \\
  \planreal{Infraestructura (Docker, MVC, MySQL)}{1}{5}{1}{5}{verdeOk}
  \planreal{APIs de usuarios y verificación}{6}{9}{6}{9}{verdeOk}
  \planreal{APIs de actividades}{10}{11}{10}{11}{verdeOk}
  \planreal{API de grupos}{12}{13}{12}{15}{ambarParcial}
  \planreal{APIs de asignaciones}{16}{18}{16}{21}{ambarParcial}
  \planreal{Integración Front--Back}{14}{15}{14}{15}{ambarParcial}
  \planreal{Revisión del contenido generado}{21}{21}{21}{21}{rojoFallo}

  \ganttvrule{Baja de un integrante}{12}
\end{ganttchart}%
}
\caption{Iteración 1: planificado (gris) frente a realidad (color).}
\label{fig:gantt-it1-real}
\end{figure}

\vspace{0.2cm}
\begin{table}[H]
\centering
\small
\rowcolors{2}{azulclaro!25}{white}
\begin{tabularx}{\textwidth}{l X X}
\rowcolor{azulPastelApagado}
\textbf{\color{white}Área} & \textbf{\color{white}Qué falló} & \textbf{\color{white}Causa principal} \\
Frontend & Implementación incompleta; la integración y la revisión no se llegaron a cerrar. & Exámenes, prácticas y entregas de otras asignaturas, más formación técnica no prevista. \\
Backend & Las APIs de grupos y asignaciones se retrasaron. & Baja de un integrante a mitad de iteración y reorganización de tareas. \\
Transversal & No hubo apoyo del backend al frontend. & El equipo de backend quedó limitado al despliegue y a asesoramiento puntual. \\
\end{tabularx}
\caption{Resumen de desviaciones de la Iteración 1.}
\end{table}

\subsection{Iteración 2: Deuda Técnica frente a Nuevo Alcance}

El siguiente diagrama coloca la Iteración 2 (12 de noviembre -- 2 de diciembre de 2026) y distingue qué trabajo es \textbf{deuda heredada de la Iteración 1} y cuál es \textbf{alcance nuevo}. Así se ve cuánto de la nueva iteración se consume en cerrar lo pendiente.

\begin{center}
\small
\cajita{ambarParcial}~Deuda técnica de la Iteración 1 \quad
\cajita{azulNuevo}~Alcance nuevo de la Iteración 2 \quad
\cajita{verdeOk}~Pruebas e integración
\end{center}

\begin{figure}[H]
\centering
\resizebox{\textwidth}{!}{%
\begin{ganttchart}[cmp, bar/.append style={draw=black!40}]{1}{21}
  \gantttitle{Noviembre 2026}{19}
  \gantttitle{Dic.}{2} \\
  \gantttitlelist{12,13,14,15,16,17,18,19,20,21,22,23,24,25,26,27,28,29,30,1,2}{1} \\

  \ganttgroup{Front-End}{1}{21} \\
  \ganttbar[bar/.append style={fill=ambarParcial}]{Inicios de sesión (profesor, estudiante) + Familia}{1}{5} \\
  \ganttbar[bar/.append style={fill=ambarParcial}]{Interfaz de crear actividad y asignar}{6}{8} \\
  \ganttbar[bar/.append style={fill=ambarParcial}]{Vista de actividad}{9}{10} \\
  \ganttbar[bar/.append style={fill=ambarParcial}]{Interfaz de agenda}{11}{14} \\
  \ganttbar[bar/.append style={fill=ambarParcial}]{Creación de usuarios y grupos}{15}{15} \\
  \ganttbar[bar/.append style={fill=ambarParcial}]{Asignación estudiante-profesor y grupos}{16}{16} \\
  \ganttbar[bar/.append style={fill=azulNuevo}]{Asignación estudiante-familiar (HU-A6)}{17}{17} \\
  \ganttbar[bar/.append style={fill=verdeOk}]{Integración Front--Back}{18}{21} \\

  \ganttgroup{Back-End}{1}{21} \\
  \ganttbar[bar/.append style={fill=azulNuevo}]{API de login familiar (HU-F1)}{1}{2} \\
  \ganttbar[bar/.append style={fill=azulNuevo}]{API de estudiantes asociados (HU-F2, HU-F5)}{3}{4} \\
  \ganttbar[bar/.append style={fill=azulNuevo}]{API de completación de actividad y pasos (HU-E7)}{5}{8} \\
  \ganttbar[bar/.append style={fill=ambarParcial}]{Apoyo al Frontend}{9}{11} \\
  \ganttbar[bar/.append style={fill=verdeOk}]{Testeo de funcionalidades}{12}{13} \\
  \ganttbar[bar/.append style={fill=verdeOk}]{Integración Front--Back}{14}{15} \\
  \ganttbar[bar/.append style={fill=ambarParcial}]{Apoyo al Frontend}{16}{17} \\
  \ganttbar[bar/.append style={fill=verdeOk}]{Integración Front--Back}{18}{19} \\
  \ganttbar[bar/.append style={fill=verdeOk}]{Testeo de funcionalidades}{20}{21} \\

  \ganttmilestone[milestone/.append style={fill=rojoFallo}]{Entrega Iteración 2}{21}
\end{ganttchart}%
}
\caption{Iteración 2: deuda técnica heredada frente a alcance nuevo.}
\label{fig:gantt-it2-deuda}
\end{figure}